\scanpage{310}
\mspage{226}
% archive sheet number: 309

prouver, sur le foncteur $f$ du cas p.~223, les deux conditions suivantes :

$\mathrm{F}1$) \deleted{\unsure{pour}} $\forall$ système \deleted{d'\unsure{objets}} de flèches \inserted{dans $M$}

\[
X(j)\xrightarrow{p(j)}\unsure{\varphi|\phi}(j)
\qquad\left|\begin{aligned}
&j\in\deleted{\unsure{I}}J_0=\mathrm{Ob}\,J\\
&\text{les }X(j)\text{ dans }\mathrm{Ob}\,C
\end{aligned}\right.
\]

il existe \inserted{un foncteur} $F:J\longrightarrow C$ et $F\xrightarrow{q}\varphi$, enfin un système de flèches $X(j)\xrightarrow{u(j)}F(j)$, \deleted{\unsure{rendant}} dans \inserted{les} $C/\varphi(j)$ i.e. rendant comm les diagrammes \unsure{i.e.,|i.e.}

\diagram{\begin{tikzcd}[ampersand replacement=\&,column sep=small]
X(j)\arrow[rr,"u(j)"]\arrow[dr,"p(j)"']\&\&F(j)\arrow[dl,"q(j)"]\\
\&\varphi(j)\&
\end{tikzcd}}

$\mathrm{F}3_{\unsure{\pi|\kappa}})$ Pour \deleted{tout} système $(F_\alpha)_{\alpha\in A}$ \deleted{\unsure{de foncteurs} $\in\underline{\mathrm{Hom}}(J,$} \inserted{$(\operatorname{card} A\leq\pi)$} d'éléments dans \deleted{les} $\underline{\mathrm{Hom}}(C,J)$, et des flèches $F_\alpha\xrightarrow{p_\alpha}\varphi$, dans $\underline{\mathrm{Hom}}(J,M)$, et pour $\deleted{\unsure{F}}G$ dans $\underline{\mathrm{Hom}}(J,\unsure{C|I})$, et $\deleted{\unsure{F}}G\xrightarrow{\deleted{\unsure{p}}q}\varphi$ dans $\underline{\mathrm{Hom}}(J,M)$, \unsure{étant donné|dont donné} un système de flèches $\unsure{u_\alpha(j)}$.

% slanted
\margin{en fait, il suffit de vérifier $\mathrm{F}3$ (cf plus bas 4.9.7.1), i.e. \unsure{lors|cas|lorsq.} $A$ réduit à un élément}

\diagram{\begin{tikzcd}[ampersand replacement=\&,column sep=small]
F_\alpha(\unsure{j})\arrow[rr,"u_\alpha(j)"]\arrow[dr,"p_\alpha(j)"']\&\&G\unsure{_\alpha|{}}(\unsure{j})\arrow[dl,"q(j)"]\\
\&\varphi(j)\arrow[ul,phantom,"\text{comm.}"{near start,swap}]\&
\end{tikzcd}}

\margin{\deleted{\unsure{système}} \deleted{\unsure{flè}} $u_\alpha(j)$ flèche dans $C/\varphi(j)$}

il existe \deleted{$\unsure{H\xrightarrow{r}\varphi}$ \unsure{dans} $\underline{\mathrm{Hom}}(J,C)_{\unsure{/\varphi}}$} \inserted{$H\xrightarrow{\unsure{q|r}}\varphi\unsure{^{\prime}|{}}$ dans $\underline{\mathrm{Hom}}(J,\unsure{C})_{/\varphi}$} et une \deleted{\unsure{un}} flèche $G\xrightarrow{v}H$ dans $\underline{\mathrm{Hom}}(J,C)_{/\varphi}$, telle que pour \unsure{tt|tout} $\alpha\inserted{{}\in A}$, le système de flèches composées

